\section{Was davon bei den Massen ankommt}
\label{sec:mass}

Die bisherige Konstruktion legt fest, wie man Zustände auswählt und unter welchen Bedingungen sie eine innere Symmetrie tragen. Das kann man schon für einen vorgegebenen Zustandsraum prüfen. Jetzt kommt eine andere Anforderung hinzu. Aus diesen Zuständen sollen konkrete Massen folgen. An den drei geladenen Leptonen können wir fragen, welche Eigenschaft eine gemeinsame Auslese haben müsste.

\subsection{Die Normbedingung im Raum der Massenwurzeln}

Für drei nichtnegative Massenwurzeln \(x_i=\sqrt{m_i}\), \(i=1,2,3\), mit \(x\ne0\) definieren wir
\begin{equation}
 K(x)=\frac{\sum_i x_i^2}{(\sum_i x_i)^2}.
 \label{eq:koide}
\end{equation}
Die physikalische Anwendung betrifft die drei geladenen Leptonen. Für die folgende Identität genügt ein beliebiger positiver Summenwert. Mit \(a=(\sum_i x_i)/3>0\) schreiben wir \(x=a(1,1,1)+d\), wobei \(\sum_i d_i=0\). Die beiden Anteile sind orthogonal. Folglich
\begin{equation}
 K(x)=\frac{3a^2+\|d\|^2}{9a^2}
      =\frac13+\frac{\|d\|^2}{9a^2},
 \qquad
 K(x)=\frac23\ \Longleftrightarrow\ \|d\|^2=3a^2 .
 \label{eq:equal-norm}
\end{equation}
Die Koide-Bedingung verlangt gleiche quadratische Normen des konstanten Anteils und der Abweichung. Ein zyklischer Basiswechsel erhält beide Normen und lässt ihr Verhältnis frei. Als Beispiel ergeben \(x=(1,2,3)\) den Wert \(7/18\) und drei gleiche Wurzeln den Wert \(1/3\). Die Rechnung verlegt die physikalische Frage damit auf die Auswahl der Normrelation \cite{physics}.

Eine für den Vergleich mit dem Spektrum nützliche Koordinate auf der
Nullsummenebene ist
\begin{equation}
 x_k=a\bigl[1+r\cos(\theta+2\pi k/3)\bigr],
 \qquad k=0,1,2,\qquad r\geq0 .
 \label{eq:koide-radius-phase}
\end{equation}
Die beiden Fourier-Vektoren mit Komponenten \(\cos(2\pi k/3)\)
und \(\sin(2\pi k/3)\) sind orthogonal und haben beide die
quadratische Norm \(3/2\). Damit gilt
\(\|d\|^2=\tfrac32a^2r^2\), und \eqref{eq:equal-norm} wird zu
\[
 K=\frac13+\frac{r^2}{6},\qquad K=\frac23\ \Longleftrightarrow\ r=\sqrt2 .
\]
Die Normrelation wählt einen Radius im Raum der normierten Wurzeln.
Entlang dieses Radius bleiben \(\theta\) und die gemeinsame
Normierung \(a\) als Modellgrößen zu bestimmen. Zusätzlich muss jeder
der drei Ausdrücke in eckigen Klammern nichtnegativ sein. Ist einer dieser Ausdrücke
negativ, liefert die Hauptwurzel der quadrierten Größe ihren
Betrag. Dann ändert sich der Nenner in \eqref{eq:koide}, und die
Radiusformel kann für dieses Spektrum nicht verwendet werden
\cite{companion,physics}.

\subsection{Der Test am Kompensator}

Auf dem euklidischen Wurzelraum sei \(e=(1,1,1)/\sqrt3\) und \(C_m=(\id+J_m)/2\) mit orthogonaler Involution \(J_m\). Angenommen, \(x=\alpha C_m e\) mit \(\alpha>0\) sei nichtnegativ und von null verschieden. Dann gilt \(\langle e,C_m e\rangle=\|C_m e\|^2\). Setzt man \(h=\langle e,J_m e\rangle\), folgt
\begin{equation}
 K(x)=\frac{1}{3\|C_m e\|^2}
      =\frac{2}{3(1+h)},\qquad
 K(x)=\frac23\ \Longleftrightarrow\ h=0.
 \label{eq:hc-criterion}
\end{equation}
Der Test verlangt also, dass \(e\) und \(J_me\) orthogonal sind. Die Projektoraxiome erlauben auch \(J_m=\id\), mit \(h=1\). Um die Orthogonalität physikalisch zu begründen, muss \(J_m\) aus dem gemeinsamen Träger folgen.

Dafür betrachtet die Physik-Ausgabe eine Koisometrie \(B:H\to\mathbb R^3\) mit \(BB^*=\id_3\) und \(J\)-invariantem Kern \cite{physics}. Da \(J\) selbstadjungiert ist, bleibt auch \((\ker B)^\perp=\operatorname{im} B^*\) invariant. \(B\) überträgt die Involution dieses orthogonalen Komplements als \(J_m=BJB^*\) auf den Wurzelraum, sodass \(BC=C_mB\). Ein unabhängig bestimmter Ausgangszustand \(f_0\) mit \(Bf_0=e\) und das physikalische Wurzelgesetz \(x=\alpha BCf_0\) würden die Bedingung auf
\[
 \langle B^*e,J B^*e\rangle=0
\]
zurückführen. Das benennt eine konkrete Prüfung am Träger. Die Auswahl von \(B\), \(f_0\) und des Wurzelgesetzes geht als eigene Voraussetzung ein. Ein rein ungerader Zustand liegt im Kern von \(C\) und liefert über diese direkte lineare Auslese keine von null verschiedenen Wurzeln.

Ein vollständig bestimmter Vergleich zeigt, welche Information in
der Wahl von \(J_m\) steckt. Wähle
\(v=(2,-1,-1)/\sqrt6\) und lasse \(J_m\) die orthonormalen
Vektoren \(e\) und \(v\) vertauschen. Auf ihrem orthogonalen
Komplement wirke \(J_m\) als Identität. Dann ergibt sich
\[
 C_m e=\frac12(e+v),\qquad
 x=\frac{\alpha}{2\sqrt3}
       \left(1+\sqrt2,\ 1-\frac1{\sqrt2},\ 1-\frac1{\sqrt2}\right).
\]
Das sind drei positive Hauptwurzeln im Vergleich, von denen zwei
gleich sind. Der Faktor vor der Klammer ist ihr
Mittelwert \(a\). Die erste Komponente hat die Richtung ausgewählt,
die in \eqref{eq:koide-radius-phase} der Phase \(\theta=0\)
entspricht. Hier ist also \(K=2/3\) bereits erreicht, bevor überhaupt
drei verschiedene Werte auftreten. Im Companion dient die
Konstruktion als Beispiel für die Realisierbarkeit von
\(\langle e,J_m e\rangle=0\). Für eine physikalische Erklärung
braucht man eine Regel, durch die der Träger gerade eine solche
Richtung und eine bestimmte Aufspaltung auswählt \cite{companion,physics}.

\subsection{Die Koide-Überlappung}

Die Überlappungsdarstellung setzt drei normierte Zustände \(v_i\) mit Gram-Matrix \(G_{ij}=\langle v_i,v_j\rangle\) an. In diesem Modell sind die Wurzeln \(x_i=s\lambda_i(G)\) mit \(s>0\). Da \(G\) positiv semidefinit ist und \(\Tr G=3\) gilt, ergibt sich
\[
 \Tr G^2=3+2\sum_{i<j}|\langle v_i,v_j\rangle|^2.
\]
Definieren wir \(\rho^2=\frac13\sum_{i<j}|\langle v_i,v_j\rangle|^2\), dann folgt
\begin{equation}
 K=\frac{\Tr G^2}{(\Tr G)^2}
   =\frac{1+2\rho^2}{3},
 \qquad K=\frac23\ \Longleftrightarrow\ \rho^2=\frac12 .
 \label{eq:overlap}
\end{equation}
Die Gleichung setzt das Wurzelgesetz für das Spektrum voraus. Sie bestimmt den Mittelwert der drei quadratischen Überlappungen. Paarweise gleiche Überlappungen sind eine zusätzliche Spezialisierung.


\subsubsection{Ein durchgehendes Beispiel mit veränderlicher Phase}

Wie kann bei festem \(K\) aus dem Beispiel mit einer doppelten
Wurzel ein Spektrum mit drei verschiedenen Wurzeln werden? Für die
im Companion untersuchte Vergleichsfamilie seien alle drei
paarweisen quadratischen Überlappungen gleich \(1/2\). Schreibe
\(U(z_1,z_2,z_3)=(z_2,z_3,z_1)\) und setze
\begin{equation}
 G(\varphi)=I_3+\frac1{\sqrt2}\left(
   e^{i\varphi/3}U+e^{-i\varphi/3}U^*\right).
 \label{eq:koide-circulant-gram}
\end{equation}
Das Produkt \(G_{12}G_{23}G_{31}=e^{i\varphi}/(2\sqrt2)\)
speichert die Bargmann-Phase \(\varphi\). Bei Umphasierung der
drei Zustände, \(v_j\mapsto e^{i\alpha_j}v_j\), wird
\(G_{ij}\) mit \(e^{i(\alpha_j-\alpha_i)}\) multipliziert.
Im Produkt heben sich diese Faktoren auf. Die Phase bleibt also
ein Merkmal des Tripels, auch wenn man die Phasen seiner drei
Vektorrepräsentanten frei wählt. Die zirkulante Darstellung
verteilt diese Phase gleichmäßig auf die drei gerichteten
Einträge \cite{physics,companion}.

Die Fourier-Basis diagonalisiert \(U\). Mit
\(\omega=e^{2\pi i/3}\) gehört zum Eigenwert \(\omega^k\) von
\(U\) die Eigenzahl
\begin{equation}
 \lambda_k(\varphi)=1+\sqrt2\cos\left(\frac{\varphi}{3}
                                     +\frac{2\pi k}{3}\right)
 \label{eq:koide-gram-spectrum}
\end{equation}
von \(G\). Für positiv semidefinites \(G\) sind dies
die in \eqref{eq:overlap} verwendeten Wurzeln in Einheiten von
\(s\). Ihr Mittelwert ist eins. Für \(\varphi=0\) ergibt sich
genau das normierte Wurzeltripel aus dem Projektorvergleich. In
der zirkulanten Familie entspricht die Änderung von \(\varphi\)
einer Änderung von \(\theta=\varphi/3\) in
\eqref{eq:koide-radius-phase}. Der Radius bleibt dabei \(\sqrt2\).
Somit kann man an diesem Beispiel direkt verfolgen, wie dieselbe
quadratische Überlappung unterschiedliche Wurzelspektren zulässt.

\subsubsection{Warum der Oktant am Rand liegt}

Für eine Realisierung durch drei Zustände muss \(G\) positiv
semidefinit sein. Die Determinantenentwicklung liefert
\begin{equation}
 \det G=1-\tfrac32+2\operatorname{Re}(G_{12}G_{23}G_{31})
       =\tfrac12(\sqrt2\cos\varphi-1).
 \label{eq:det}
\end{equation}
Alle \(2\times2\)-Hauptminoren sind gleich \(1/2>0\).
Im Hauptwertbereich der Phase ist \(G\) daher für
\(|\varphi|<\pi/4\) positiv definit. Bei \(|\varphi|=\pi/4\)
hat die Matrix Rang zwei, und für \(|\varphi|>\pi/4\) erhält
man keine Gram-Matrix von Zustandsvektoren. Die Determinante
zeigt damit, an welcher Stelle der formale Ausdruck für das
Spektrum in \eqref{eq:koide-gram-spectrum} seine Interpretation
als drei nichtnegative Hauptwurzeln verliert.

\begin{figure}[!htbp]
\centering
\begin{tikzpicture}[>=Latex,x=9.6cm,y=4.25cm,font=\small]
  \fill[qs blue!10] (-.25,0) rectangle (.25,.255);
  \draw[qs axis,->] (-.57,0)--(.58,0) node[right] {\(\varphi/\pi\)};
  \draw[qs axis,->] (0,-.39)--(0,.30) node[above] {\(\det G\)};
  \foreach \x/\lab in {-.5/{-\frac12},-.25/{-\frac14},.25/{\frac14},.5/{\frac12}} {
    \draw[qs axis] (\x,-.012)--(\x,.012);
    \node[below=2pt] at (\x,0) {\(\lab\)};
  }
  \draw[very thick,qs blue,domain=-.45:.45,samples=81]
     plot (\x,{(sqrt(2)*cos(180*\x)-1)/2});
  \draw[densely dashed,qs axis] (-.25,0)--(-.25,.255) (.25,0)--(.25,.255);
  \fill[qs blue] (-.25,0) circle (1.65pt) (.25,0) circle (1.65pt);
  \node[qs ink,fill=qs blue!10,inner sep=2pt] at (0,.09) {positiv definit};
  \node[qs ink,anchor=west] at (.42,-.18) {\(\det G<0\)};
\end{tikzpicture}
\caption{Determinante der Gram-Matrix bei festen paarweisen quadratischen Überlappungen \(1/2\). Im gezeigten Phasenintervall entsprechen die beiden Nullstellen dem Rangabfall. Im markierten Bereich zwischen ihnen sind alle drei Eigenwerte positiv. In der gesamten zulässigen Familie hat die Koide-Auslese denselben Wert.}
\label{fig:koide-phase}
\end{figure}


Bei \(\varphi=\pi/4\) erhält man
\[
 \left(\lambda_0,\lambda_1,\lambda_2\right)
  =\left(\frac{3+\sqrt3}{2},\,0,\,\frac{3-\sqrt3}{2}\right).
\]
Den Rand kann man direkt in \(\mathbb C^2\) realisieren. Wähle
\(v_1=(1,0)\), \(v_2=(1,1)/\sqrt2\) und
\(v_3=(1,i)/\sqrt2\). Die drei Überlappungen haben den
Betrag \(1/\sqrt2\) und ihr gerichtetes Produkt ist
\((1+i)/4\). Es sind die drei Spin-\(\tfrac12\)-Zustände
eines Bloch-Oktanten, den die Physik-Ausgabe diskutiert.
Die positive Matrix im Inneren des Phasenintervalls hat
hingegen Rang drei und verlangt eine Realisierung in
mindestens drei komplexen Dimensionen. Eine mögliche
Realisierung liefern dann die drei Spalten von \(G^{1/2}\).
Für \(0<|\varphi|<\pi/4\) sind alle drei Wurzeln positiv
und verschieden. Der Wert in \eqref{eq:overlap} ist an
Rand und Innerem derselbe. Die Rechnung legt daher die
zulässige Phasenfamilie offen, während die Physik erklären
müsste, welcher Wert in dieser Familie tatsächlich
ausgewählt wird. Die Interpretation des Defizits
\(\pi/4-|\varphi|\) als Kern in Abschnitt~5.8 der
Physik-Ausgabe bleibt eine bedingte Konstruktion
\cite{physics,companion}.

\subsubsection{Welches Vorzeichen der Oktantschritt hat}

Der Oktant legt den Betrag der Phase fest, nicht ihr Vorzeichen. Durchläuft
man dieselben drei Zustände in umgekehrter Reihenfolge,
\(v_1\to v_3\to v_2\to v_1\), bleiben Anfangs- und Endstrahl und alle drei
Übergangswahrscheinlichkeiten \(\tfrac12\) gleich. Das gerichtete Produkt
wird aber zu \((1-i)/4\), der Schritt also zu \(-\pi/4\) statt \(+\pi/4\).
Beträge und Endpunkte entscheiden das Vorzeichen daher nicht. Auch eine
gleichmäßige Verteilung der aufsummierten Phasen sagt nichts: Die beiden
konstanten Schrittfolgen \(\pm1\) modulo acht durchlaufen in einem Zyklus
alle acht Phasen \cite{h116}.

Der Kompensator liefert das Vorzeichen ebenfalls nicht. Auf den
geschlossenen ungeraden Generationszuständen gilt \(JS=-S\), also
\(CS=(S+JS)/2=0\). Eine reelle Projektion vom Rang fünf, die eine reelle
Richtung der komplexen Nullgeraden der Oktant-Auslese entfernt, gibt es
zwar. Die beiden reellen Richtungen \(n\) und \(in\) dieser Geraden liefern
aber dieselben Ränge und dieselbe Auslese, und \(J\) wählt nicht zwischen
ihnen \cite{h115}.

Übrig bleibt die Auswahl des geordneten Rückwegs. Ein Modell mit zwei
Fenstern macht den Zeitpunkt sichtbar. Das erste Fenster hält die
zulässigen Orientierungen als Maske in einem vorläufigen Zustand, das
zweite prüft sie mit einer zweiten Maske, und ein Vorzeichen entsteht erst,
wenn genau eine Orientierung übrig bleibt. Von den sechzehn Maskenpaaren
ergeben drei \(+\pi/4\), drei \(-\pi/4\), neun sind gesperrt und eines
bleibt mehrdeutig. Die vorhandenen Regeln des Companion vertragen sich mit
diesem Zeitpunkt, enthalten aber keine Größe, die zwischen den beiden
Vorzeichen entscheidet \cite{h117}. Alle drei Rechnungen waren vorab
festgelegt, ihr Ausgang stand aber analytisch fest; sie gelten deshalb
nicht als blinde Tests. Die Physik-Ausgabe von v5.3.0 fasst sie in Abschnitt~5.9 zusammen \cite{physics}.
Eine spätere Prüfung ordnet die Orientierung ein: Frame-Abbildungen und
Parität erhalten sie, die Zeitumkehr der Feldamplituden wirkt als komplexe
Konjugation und kehrt sie um. Das Vorzeichen ist also eine Zeitorientierung
der Auslese, und keine im Companion festgehaltene Struktur legt es fest
\cite{h118}.

\subsubsection{Was sich beim Kombinieren von Präparationen ändert}

Die Phase kann nur dann als Eigenschaft einer Auslese
verwendet werden, wenn feststeht, an welcher Stelle der
Auslese sie eingeht. Betrachten wir zwei der oben
beschriebenen positiven Tripel mit \(\theta_\pm=\theta_0\pm h\)
im selben positiv zulässigen Intervall, demselben \(a\)
und derselben Zuordnung der Komponenten in einem
festen Materieframe. Wenn die beiden Wurzeln einer
Komponente zuerst zu gleichen Teilen gemittelt werden,
folgt mit der Addition der beiden Kosinuswerte
\begin{equation}
 \bar x_k=a\bigl[1+\sqrt2\cos h\,
                      \cos(\theta_0+2\pi k/3)\bigr],\qquad
 K(\bar x)=\frac{1+\cos^2 h}{3}.
 \label{eq:koide-mixture}
\end{equation}
Die Präparationen dürfen also einzeln \(K=2/3\) erfüllen,
während die gemittelte Wurzelauslese für eine echte
Phasendifferenz einen kleineren Radius hat. Würde man die
beiden Massen \(x_{+k}^2,x_{-k}^2\) unmittelbar mitteln,
ergäbe sich pro Komponente
\[
 \frac{x_{+k}^2+x_{-k}^2}{2}
   =\bar x_k^{\,2}+\frac{(x_{+k}-x_{-k})^2}{4}.
\]
Das zweite Ergebnis enthält einen Varianzterm. Beide
Verfahren gehen von denselben Präparationen aus und
unterscheiden sich in ihrer Operationsreihenfolge.

Noch einmal eine eigene Möglichkeit ist die kohärente
Addition zweier Pfadbeiträge \(B_1+e^{i\phi}B_2\) zu
einem Übergangsoperator \(B\). Im Produkt \(B^*B\)
entstehen dann die Kreuzterme
\(e^{i\phi}B_1^*B_2+e^{-i\phi}B_2^*B_1\).
Eine relative Phase kann deshalb im Spektrum des
Operators wirken. Ein gemeinsamer Phasenfaktor vor
\(B\) fällt aus \(B^*B\) heraus. Für eine native
Zustandspräparation müsste man angeben, welche dieser
Operationen physikalisch realisiert wird. Die
Phasenmittelung ist im Companion ein ausdrücklich
bedingter Vergleich \cite{companion,physics}.

\subsection{Eine Vergleichsdynamik mit leichter und schwerer Komponente}

Die Kompression aus Abschnitt \ref{sec:matter} kann mit einem Übergangsoperator zu einem zusätzlichen Raum verbunden werden. Als konkreten Vergleich betrachten wir einen Übergang \(B_{\mathrm{lh}}\) zwischen zwei endlichdimensionalen Hilberträumen mit ihren deklarierten positiven Metriken. Der schwere Block habe die konstante Lücke \(\Delta>0\) und die Kopplung das Gewicht \(\mu>0\). In orthonormalen Basen lautet der hermitesche Blockoperator
\[
 \Omega=
 \begin{pmatrix}
 0 & \mu B_{\mathrm{lh}}^* \\
 \mu B_{\mathrm{lh}} & \Delta\id
 \end{pmatrix}.
\]
Die im Companion behandelte Form mit gewichteten Metriken geht durch einen Metrikwechsel in diese Form über \cite{companion,repo}. Gehört \(\sigma_i>0\) zu den Singulärwerten des Übergangs, reduziert sich der zugehörige Teil auf die Matrix mit Diagonale \(0,\Delta\) und Nebendiagonale \(\mu\sigma_i\). Die negative Eigenzahl \(-m_i\) liefert
\begin{equation}
 m_i(m_i+\Delta)=\mu^2\sigma_i^2 .
 \label{eq:gap}
\end{equation}
Für den Vergleich sind \(\sigma_i\) dimensionslos und \(\mu,\Delta,m_i\) haben die Dimension einer Masse. Für \(\mu\sigma_i\ll\Delta\) ist \(m_i\) in führender Ordnung \(\mu^2\sigma_i^2/\Delta\). Diese Näherung hängt von der Lücke ab. Der leichte Eigenzustand hat bei von null verschiedener Kopplung auch eine Komponente im schweren Block.

Das zugehörige Eigenwertproblem lässt sich mit nur zwei
Komponenten verfolgen. Seien \(u_i,w_i\) normierte
Singularvektoren mit \(B_{\mathrm{lh}}u_i=\sigma_iw_i\)
und \(B_{\mathrm{lh}}^*w_i=\sigma_iu_i\). Für den Eigenzustand
zur Eigenzahl \(-m_i\) sind die Koeffizienten in dieser Basis
proportional zu
\[
 \left(1,-\frac{\mu\sigma_i}{\Delta+m_i}\right).
\]
Damit ist der quadrierte Anteil der schweren Komponente
\(m_i/(\Delta+2m_i)\), nach Einsetzen von \eqref{eq:gap}
in die Normierung. Ein frei gewählter Übergang liefere etwa
\(q_i=\sigma_i^2\). Mit der dimensionslosen Kopplung
\(\eta=\mu/\Delta\) folgt exakt
\begin{equation}
 m_i=\frac{\Delta}{2}\left(\sqrt{1+4\eta^2q_i}-1\right).
 \label{eq:koide-scale-gap}
\end{equation}
Bei festem \(q_i,\eta\) multipliziert eine Änderung von
\(\Delta\) alle \(m_i\) mit demselben Faktor. Ändert man
\(\eta\) bei festem \(q_i\), ändern sich im Allgemeinen auch
die Massenverhältnisse. Der Parameter für
die relative Wirkung der Kopplung und die absolute
Massenskala haben damit getrennte Aufgaben. Der Anteil
der schweren Komponente im Eigenzustand hängt von der
dimensionslosen Wahl ab und bleibt bei einer gemeinsamen
Skalierung von Lücke und Kopplung gleich
\cite{companion,repo}.

Die rechte Seite von \eqref{eq:gap} setzt den Übergangsoperator voraus. Die Gram-Auslese in \eqref{eq:overlap} setzt dagegen ein Gesetz für die \emph{Wurzeln} der Massen voraus. Eine Identifikation der Ergebnisse würde deshalb eine Abbildung benötigen, die aus dem Übergang und seinen Skalen genau das Wurzelspektrum erzeugt. So wird die native Operatoridentifikation zu einer konkreten Anforderung an zwei explizite Modelle.

\subsection{Was ein Massengesetz zusätzlich festlegen muss}

Im Gram-Beispiel stehen damit drei unterschiedliche Entscheidungen
an. Die Wahl der Überlappungen fixiert die Normrelation. Die Phase
bestimmt innerhalb des gezeigten Intervalls, welches Verhältnis
der Wurzeln erreicht wird. Der Faktor \(s\) gibt den Wurzeln ihre
Einheit und setzt die absolute Massenskala über \(s^2\). Der
Projektorvergleich liefert ein Kriterium für die Normrelation.
Die Blockdynamik liefert ein Beispiel dafür, wie in einer
anderen Auslese ein dimensionaler Parameter eingeht. Um diese
Stücke zu verbinden, müsste der gemeinsame Träger die
Zustandsauswahl aus Abschnitt \ref{sec:matter} und einen
nativen Operator liefern, dessen Wurzelspektrum gerade diese
Eigenschaften besitzt. Die Auswahl muss für einen Test vor dem
Vergleich mit den Zielmassen feststehen.

Im folgenden Abschnitt geht es um die nächste Stufe dieser
Frage. Ein vorgegebener Massenterm lässt sich in eine
dynamische Gleichung einsetzen. Wie der native Massenoperator
diesen Term liefern und wie derselbe Materiezustand in einer
geometrischen Quelle auftreten würde, wird damit zu einer
konkreten Anforderung an die Kopplung.
